\documentclass{article}
\usepackage{amsmath}
\usepackage{amssymb}

\title{Comprehensive rtex PDF Test}
\author{rtex Test Suite}
\date{\today}

\begin{document}
\maketitle

\begin{abstract}
This document exercises the major LaTeX features supported by rtex in a single
conversion pass: metadata, sections, text formatting, mathematics, lists, tables,
theorem environments, links, colors, footnotes, and code blocks.
\end{abstract}

\section{Document Structure}
\subsection{Headings and Text}
Plain paragraphs with \textbf{bold}, \textit{italic}, \emph{emphasis},
\texttt{monospace}, \textsc{Small Caps}, and \underline{underlined} text.
Superscripts like $x\textsuperscript{2}$ and subscripts like $H\textsubscript{2}O$
appear inline.

\section*{Unnumbered Section}
Starred sections omit numbers but still render body text.

\section{Mathematics}

\subsection{Inline and Display}
Inline: $E = mc^2$, $\alpha + \beta = \gamma$, and $\int_0^1 x^2\,dx = \tfrac{1}{3}$.
Display:
\begin{equation}
    x = \frac{-b \pm \sqrt{b^2 - 4ac}}{2a}
\end{equation}

\subsection{Aligned and Gathered Equations}
\begin{align}
    a &= b + c \\
    d + e &= f + g + h \\
    i &= j
\end{align}

\begin{gather}
    x = 1 \\
    y = 2 \\
    z = 3
\end{gather}

\begin{multline}
    a + b + c + d + e + f + g + h + i + j \\
    + k + l + m
\end{multline}

\subsection{Piecewise Functions and Cases}
\[
    f(x) = \begin{cases}
        x^2  & \text{if } x \geq 0 \\
        -x   & \text{if } x < 0
    \end{cases}
\]

Binomial coefficients: $\binom{n}{k}$ and ellipses $\ldots$, $\cdots$, $\vdots$, $\ddots$.

\section{Lists}

\subsection{Unordered and Ordered}
\begin{itemize}
    \item First item with math $\pi \approx 3.14$
    \item Second item
    \item Third item
\end{itemize}

\begin{enumerate}
    \item Step one
    \item Step two
    \item Step three
\end{enumerate}

\subsection{Nested Lists}
\begin{itemize}
    \item Outer alpha
    \begin{itemize}
        \item Inner one
        \item Inner two
    \end{itemize}
    \item Outer beta
\end{itemize}

\subsection{Description List}
\begin{description}
    \item[Parser] Converts TeX source into a structured AST.
    \item[Builder] Walks the AST and emits PDF content streams.
    \item[Formatter] Maps LaTeX math commands to Unicode symbols.
\end{description}

\section{Tables}
\begin{tabular}{lcc}
\hline
Quantity & Symbol & Value \\
\hline
Speed of light & $c$ & $3 \times 10^8$ m/s \\
Planck constant & $h$ & $6.626 \times 10^{-34}$ J$\cdot$s \\
Fine-structure & $\alpha$ & $\approx 1/137$ \\
\hline
\end{tabular}

\section{Theorem Environments}
\begin{theorem}[Pythagoras]
For a right triangle with legs $a, b$ and hypotenuse $c$, $a^2 + b^2 = c^2$.
\end{theorem}

\begin{proof}
By similarity of triangles. \qed
\end{proof}

\begin{definition}
A \emph{group} is a set $G$ with an associative binary operation and identity.
\end{definition}

\begin{lemma}
Every finite group has a well-defined order.
\end{lemma}

\section{Links, Colors, and Notes}
Read the docs at \href{https://example.com/rtex}{rtex documentation} or visit
\url{https://example.com}.

Colored text: \textcolor{red}{Important warning} and \textcolor{blue}{Information}.

A footnote marker\footnote{This footnote verifies bottom-of-page rendering.} ends the paragraph.

\section{Code and Quotes}
\begin{quote}
    ``Programs must be written for people to read, and only incidentally for
    machines to execute.'' --- Abelson \& Sussman
\end{quote}

\begin{verbatim}
fn main() {
    println!("Hello from rtex");
}
\end{verbatim}

\section{References and Bibliography}
See \cite{knuth1984} and \cite{lamport1994} for background. Equation~\eqref{eq:sample}
uses a label below.

\begin{equation}
    \label{eq:sample}
    \nabla \cdot \mathbf{E} = \frac{\rho}{\varepsilon_0}
\end{equation}

Cross-reference: Section~\ref{sec:conclusion}.

\begin{thebibliography}{9}
\bibitem{knuth1984} D.~E. Knuth, \emph{The TeXbook}, Addison-Wesley, 1984.
\bibitem{lamport1994} L.~Lamport, \emph{LaTeX: A Document Preparation System}, 2nd ed., 1994.
\end{thebibliography}

\section{Conclusion}
\label{sec:conclusion}
If this document converts to a valid multi-section PDF, rtex covers the core
LaTeX authoring surface area end to end.

\end{document}
